\documentclass[../../../uem_phy_std_a.tex]{subfiles}

\begin{document}
    \begin{enumerate}
        \item Yについての電荷保存則より，
            
            \myspaceb%
            $\displaystyle%
                -x_1 + y_1 = 0 \;.
            $\\
            キルヒホッフ則より，
            
            \myspaceb%
            $\displaystyle%
                \myfracc{y_1}{\varepsilon_0 S} \cdot 2d%
                + \myfracc{x_1}{\varepsilon_0 S} \cdot d = 2V_0%
                \qquad \therefore \; x_1 = y_1%
                = \uwave{\myfracc{2\varepsilon_0 S}{3d} V_0} \;.
            $
        \item 電荷保存則より，Xの電荷は変わらない．つまり，
            
            \myspaceb%
            $\displaystyle%
                x_2 = x_1 = \uwave{\myfracc{2\varepsilon_0 S}{3d} V_0} \;.
            $\\
            キルヒホッフ則より，
            
            \myspaceb%
            $\displaystyle%
                \myfracc{y_2}{\varepsilon_0 S} \cdot 2d = V_0%
                \qquad \therefore \; y_2 = \uwave{\myfracc{\varepsilon_0 S}{2d} V_0} \;.
            $
        \item Yについての電荷保存則より，
            
            \myspaceb%
            $\displaystyle%
                -x_3 + y_3 = -x_2 + y_2%
                \qquad \therefore \; -x_3 + y_3 = -\myfracc{\varepsilon_0 S}{6d}V_0 \;.
            $\\
            キルヒホッフ則より，
            
            \myspaceb%
            $\displaystyle%
                \myfracc{y_3}{\varepsilon_0 S} \cdot 2d%
                + \myfracc{x_3}{\varepsilon_0 S} \cdot d = 2V_0
            $

            \myspaceb%
            $\displaystyle%
                \therefore \; x_3 = \uwave{\myfracc{7\varepsilon_0 S}{9d} V_0} \;,%
                \quad y_3 = \uwave{\myfracc{11\varepsilon_0 S}{18d} V_0} \;.
            $
        \item Yについての電荷保存則より，
            
            \myspaceb%
            $\displaystyle%
                -x_4 + y_4 = -x_3 + y_3%
                \qquad \therefore \; -x_4 + y_4 = -\myfracc{\varepsilon_0 S}{6d}V_0 \;.
            $\\
            キルヒホッフ則より，
            
            \myspaceb%
            $\displaystyle%
                \myfracc{y_4}{\varepsilon_0 S} \cdot 2d%
                + \myfracc{x_4}{\varepsilon_0 \varepsilon_r S} \cdot d = 2V_0
            $

            \myspaceb%
            $\displaystyle%
                \therefore \;%
                \begin{myaligna}[t]
                    & x_4 = \uwave{\myfracc{7\varepsilon_r}{3(2\varepsilon_r + 1)}%
                    \myfracc{\varepsilon_0 S}{d} V_0} \;,\\
                    & y_4 = \uwave{\myfracc{12\varepsilon_r - 1}{6(2\varepsilon_r + 1)}%
                    \myfracc{\varepsilon_0 S}{d} V_0} \;.
                \end{myaligna}
            $
    \end{enumerate}
\end{document}
